\documentclass[11pt]{article}

\usepackage[margin=1in]{geometry}
\usepackage{amsmath,amssymb,amsthm,mathtools,bm}
\usepackage{booktabs,enumitem,microtype,xcolor}
\usepackage[hidelinks]{hyperref}
\usepackage[capitalize,noabbrev]{cleveref}

\definecolor{established}{HTML}{166534}
\definecolor{proposed}{HTML}{1D4ED8}
\definecolor{open}{HTML}{9A3412}

\newtheorem{definition}{Definition}[section]
\newtheorem{proposition}[definition]{Proposition}
\newtheorem{lemma}[definition]{Lemma}
\newtheorem{remark}[definition]{Remark}
\newtheorem{question}[definition]{Research question}
\newtheorem{milestone}[definition]{Milestone}

\newcommand{\R}{\mathbb{R}}
\newcommand{\E}{\mathbb{E}}
\newcommand{\Prob}{\mathbb{P}}
\newcommand{\op}{\mathrm{op}}
\newcommand{\tr}{\mathrm{tr}}
\newcommand{\cH}{\mathcal{H}}
\newcommand{\cE}{\mathcal{E}}
\newcommand{\norm}[1]{\left\lVert #1\right\rVert}
\newcommand{\ip}[2]{\left\langle #1,#2\right\rangle}
\newcommand{\status}[2]{\textcolor{#1}{\textbf{#2}}}

\title{Robust Kernel Ridge Regression via Random Fourier Features\\
\large A Master's-Level Research Project Prospectus}
\author{}
\date{September 2026}

\begin{document}
\maketitle

\begin{abstract}
This project studies kernel ridge regression when an adversary may replace an
$\epsilon$ fraction of the training pairs.  The proposed route first replaces a
shift-invariant kernel by a finite random Fourier feature map and then applies
robust finite-dimensional estimation.  Random Fourier features are uniformly
bounded, which removes some tail difficulties, but boundedness alone does
\emph{not} supply the relative moment conditions used by many sharp robust
covariance results.  Moreover, estimating a covariance matrix in operator norm
does not justify whitening directions whose population variance lies below the
estimation error.  The project therefore asks whether ridge regularization,
spectral truncation, and a source condition can make the unresolved directions
irrelevant for prediction.  This prospectus separates established ingredients
from proposed claims, identifies the main mathematical gap, and gives a staged
plan suitable for a master's student.
\end{abstract}

\tableofcontents

\section{Project overview}

Kernel ridge regression (KRR) combines a flexible nonlinear model with a
convex optimization problem.  Its direct implementation, however, uses an
$n\times n$ kernel matrix and is vulnerable to a small number of maliciously
chosen training examples.  Random Fourier features (RFFs) address the first
issue by turning a shift-invariant kernel into an explicit $m$-dimensional
feature map.  This project asks whether the same reduction can also help with
the second issue by converting robust KRR into robust ridge regression.

The intended pipeline is
\[
 \text{KRR}\longrightarrow \text{RFF ridge regression}
 \longrightarrow \text{robust moment estimation}
 \longrightarrow \text{regularized or truncated solution}.
\]
The central obstacle is spectral.  The covariance of the RFF vector can have
very small eigenvalues.  An additive operator-norm estimate cannot resolve
these directions, and naively inverting the estimated covariance amplifies the
error.  Prediction loss, on the other hand, downweights low-variance
directions.  This suggests that global conditioning may be unnecessary if the
target has little prediction-relevant mass in the unresolved spectrum.

\paragraph{Primary objective.}
Prove and test a finite-dimensional result for robust ridge regression on RFF
features that controls clean-population prediction error without assuming a
uniform lower bound on the feature covariance.

\paragraph{Secondary objective.}
Determine when a standard source condition for the kernel integral operator
implies the finite-dimensional spectral alignment needed by the RFF theorem.

\paragraph{Status convention.}
Throughout the document,
\status{established}{Established} denotes a standard fact or a claim proved in
the text; \status{proposed}{Proposed} denotes a plausible target whose precise
assumptions and rate remain to be proved; and \status{open}{Open} denotes a
genuine technical question.

\section{Statistical model and performance criterion}

Let $P$ be a distribution on $\mathcal X\times\R$.  A clean pair $(X,Y)\sim P$
satisfies
\begin{equation}\label{eq:model}
  Y=f_\star(X)+\xi, \qquad \E[\xi\mid X]=0.
\end{equation}
The noise assumptions will be specified for each theorem; useful starting
points are bounded or conditionally sub-Gaussian noise.  Later work may relax
these to finite conditional moments using robust estimation of the
feature--response cross-moment.

\begin{definition}[Strong replacement contamination]
Draw an i.i.d. clean sample $S=((X_i,Y_i))_{i=1}^n\sim P^n$.  After observing
$S$ and the realized random feature map, an adversary selects a set
$O\subseteq[n]$ with $|O|\leq \epsilon n$ and replaces every pair indexed by
$O$ arbitrarily.  The learner observes only the resulting sample
$\widetilde S$.
\end{definition}

This is stronger than a model in which only labels are corrupted.  It also
makes explicit that the RFF randomness is not hidden from the adversary.

The main loss is clean-population prediction error:
\begin{equation}\label{eq:prediction}
 \norm{f-f_\star}_{L^2(P_X)}^2
 =\E_{X\sim P_X}\bigl[(f(X)-f_\star(X))^2\bigr].
\end{equation}
Coefficient error is deliberately secondary: it can be large in a direction
that has negligible variance and hence negligible predictive importance.

\section{Clean kernel ridge regression}

Let $k$ be a bounded positive-semidefinite kernel with reproducing kernel
Hilbert space (RKHS) $\cH_k$.  Clean population ridge regression solves
\begin{equation}\label{eq:population-krr}
 f_\lambda=\arg\min_{f\in\cH_k}
 \left\{\E(Y-f(X))^2+\lambda\norm{f}_{\cH_k}^2\right\},
\end{equation}
while empirical KRR replaces the expectation by the sample average.  By the
representer theorem, the empirical minimizer has the form
$\widehat f_\lambda(\cdot)=\sum_{i=1}^n\alpha_i k(X_i,\cdot)$, where, under the
normalization in \cref{eq:population-krr},
\[
  \bm\alpha=(K+n\lambda I_n)^{-1}\bm y.
\]

The population kernel integral operator $T:L^2(P_X)\to L^2(P_X)$ is
\begin{equation}\label{eq:T}
  (Tg)(x)=\int k(x,x')g(x')\,dP_X(x').
\end{equation}
Two standard quantities governing clean KRR are the effective dimension
\begin{equation}\label{eq:effective-dim}
  \mathcal N(\lambda)=\tr\bigl(T(T+\lambda I)^{-1}\bigr)
\end{equation}
and a source condition such as
\begin{equation}\label{eq:kernel-source}
  f_\star=T^r g,\qquad \norm{g}_{L^2(P_X)}\leq R.
\end{equation}
The exact admissible range and interpretation of $r$ depend on the convention
and ambient space.  The student should not import a rate from the literature
without first matching these definitions.  Caponnetto and De Vito
\cite{caponnetto2007} and the correction by Sutherland \cite{sutherland2017}
are useful references for this operator viewpoint.

\section{Random Fourier features}

Suppose $\mathcal X\subseteq\R^d$ and $k(x,x')=\kappa(x-x')$ is continuous,
shift invariant, and normalized so that $k(x,x)=1$.  Bochner's theorem implies
that $\kappa$ is the Fourier transform of a probability measure $\rho$:
\[
 \kappa(t)=\int_{\R^d} e^{i\omega^\top t}\,d\rho(\omega).
\]
For the Gaussian kernel
\[
 k_\sigma(x,x')=\exp\!\left(-\frac{\norm{x-x'}_2^2}{2\sigma^2}\right),
\]
$\rho$ is Gaussian with covariance $\sigma^{-2}I_d$.  Draw
$\omega_1,\ldots,\omega_m\stackrel{\mathrm{iid}}\sim\rho$ and
$b_1,\ldots,b_m\stackrel{\mathrm{iid}}\sim\operatorname{Unif}[0,2\pi]$, and
define
\begin{equation}\label{eq:rff}
 \Phi(x)=\sqrt{\frac{2}{m}}
 \bigl(\cos(\omega_1^\top x+b_1),\ldots,
       \cos(\omega_m^\top x+b_m)\bigr)^\top.
\end{equation}
Then $\E_{\omega,b}\ip{\Phi(x)}{\Phi(x')}=k(x,x')$; see Rahimi and Recht
\cite{rahimi2007}.  Ridge regression in this feature space is
\begin{equation}\label{eq:rff-ridge}
 \widehat\beta_\lambda=\arg\min_{\beta\in\R^m}
 \left\{\frac1n\sum_{i=1}^n(Y_i-\beta^\top\Phi(X_i))^2
 +\lambda\norm{\beta}_2^2\right\},
 \qquad \widehat f(x)=\widehat\beta_\lambda^\top\Phi(x).
\end{equation}

\subsection{What bounded features do and do not imply}

Let $Z=\Phi(X)$.  The chosen normalization gives, deterministically,
\begin{equation}\label{eq:bounded}
 \norm{Z}_2^2=\frac{2}{m}\sum_{j=1}^m
 \cos^2(\omega_j^\top X+b_j)\leq 2.
\end{equation}
Consequently $|v^\top Z|\leq\sqrt2\norm v_2$ and
$\E|v^\top Z|^4\leq4\norm v_2^4$ for every $v$.

This is useful absolute control, but the following distinction is essential.
Sharp covariance results under weak moments often assume a \emph{relative}
marginal moment condition
\begin{equation}\label{eq:l4l2}
 \bigl(\E|v^\top(Z-\E Z)|^4\bigr)^{1/4}
 \leq L\bigl(\E|v^\top(Z-\E Z)|^2\bigr)^{1/2}
 \quad\text{for all }v.
\end{equation}
Equation \eqref{eq:bounded} does not imply \eqref{eq:l4l2} with a uniform $L$:
the variance in the denominator may be arbitrarily small.  Thus the statement
``RFFs are bounded, therefore a relative operator-norm covariance theorem
applies'' is not valid without an additional argument or assumption.

\begin{remark}[Centering and second moments]
For regression, the normal matrix is usually the second moment
$M=\E[ZZ^\top]$, not the centered covariance
$\Sigma=\E[(Z-\E Z)(Z-\E Z)^\top]$.  They agree only when $\E Z=0$.
One may include an unpenalized intercept and center the features robustly, or
work directly with $M$.  This prospectus uses $M$ unless centering is explicit.
\end{remark}

\section{The finite-dimensional ridge problem}

Condition on the realized feature map and define
\begin{equation}\label{eq:moments}
 M=\E[ZZ^\top],\qquad c=\E[ZY],\qquad
 \beta_\lambda=(M+\lambda I)^{-1}c.
\end{equation}
The corresponding population RFF predictor is
$f_{m,\lambda}(x)=\beta_\lambda^\top\Phi(x)$.  If robust estimators
$\widehat M$ and $\widehat c$ are available, the plug-in estimator is
\begin{equation}\label{eq:plugin}
 \widetilde\beta_\lambda=(\widehat M+\lambda I)^{-1}\widehat c.
\end{equation}

The ridge shift already prevents literal inversion of tiny eigenvalues.
Nevertheless, the basic resolvent identity shows where a potentially harmful
$1/\lambda$ factor enters:
\begin{align}
 \widetilde\beta_\lambda-\beta_\lambda
 ={}&(\widehat M+\lambda I)^{-1}(\widehat c-c)\nonumber\\
 &+(\widehat M+\lambda I)^{-1}(M-\widehat M)\beta_\lambda.
 \label{eq:resolvent}
\end{align}
An absolute bound $\norm{\widehat M-M}_\op\leq\delta$ combined only with
$\norm{(\widehat M+\lambda I)^{-1}}_\op\leq1/\lambda$ yields a term of order
$\delta/\lambda$.  The project aims to replace this worst-case Euclidean
argument by one adapted to the prediction geometry.

For any coefficient error $h$,
\begin{equation}\label{eq:pred-norm}
 \norm{h^\top\Phi(\cdot)}_{L^2(P_X)}^2=h^\top M h
 =\norm{M^{1/2}h}_2^2.
\end{equation}
Thus the natural perturbation object is not merely
$\norm{\widehat M-M}_\op$, but the ridge-relative error
\begin{equation}\label{eq:relative}
 \Delta_\lambda=
 \norm{(M+\lambda I)^{-1/2}(\widehat M-M)
       (M+\lambda I)^{-1/2}}_\op.
\end{equation}
When $\Delta_\lambda<1$, standard matrix perturbation arguments compare the
true and estimated ridge geometries without requiring $M$ to be invertible.
Obtaining such a guarantee robustly is one of the main technical tasks.

\section{Spectral resolution and source alignment}

Let $M=\sum_{j=1}^m\mu_j u_ju_j^\top$, with $\mu_j\geq0$.  Suppose only that
$\norm{\widehat M-M}_\op\leq\delta$.  Eigenvalues far below $\delta$ and
eigenspaces separated by gaps comparable to $\delta$ cannot be expected to be
accurately recovered.  Full whitening by $\widehat M^{-1/2}$ is therefore
unstable.  Ridge regularization or thresholding at a level
$\tau\gtrsim\delta$ is the appropriate default.

Consider the finite-dimensional source condition
\begin{equation}\label{eq:finite-source}
  \beta=M^r v,\qquad \norm v_2\leq R,\qquad r\geq0.
\end{equation}
This says that the target coefficient decreases along low-variance
eigendirections.  Let $P_{<\tau}=\sum_{\mu_j<\tau}u_ju_j^\top$.

\begin{proposition}[Prediction cost of discarding the bottom spectrum]
Under \eqref{eq:finite-source},
\begin{equation}\label{eq:trunc-bound}
 \norm{M^{1/2}P_{<\tau}\beta}_2
 \leq R\tau^{r+1/2}.
\end{equation}
\end{proposition}
\begin{proof}
Write $v=\sum_jv_ju_j$.  Then
\[
 \norm{M^{1/2}P_{<\tau}\beta}_2^2
 =\sum_{\mu_j<\tau}\mu_j^{2r+1}v_j^2
 \leq\tau^{2r+1}\norm v_2^2.
\]
Taking square roots proves the claim.
\end{proof}

This elementary inequality is the main motivation for the project, but it is
not yet a robust regression theorem.  Such a theorem must also control the
estimated top spectral subspace, the cross-moment $c$, sampling error, and the
interaction between $\lambda$ and $\tau$.

\section{Robust estimation ingredients}

\subsection{Estimating the feature second moment}

The matrices $Z_iZ_i^\top$ satisfy
$\norm{Z_iZ_i^\top}_\op\leq2$ for clean observations.  This boundedness gives
standard concentration for the clean empirical second moment.  Under strong
contamination, however, the corrupted feature vectors need not be bounded: the
adversary replaces the observed pairs and is not required to provide points
that are consistent with $\Phi(x)$.  A robust procedure must therefore reject,
filter, truncate, or reweight arbitrary submitted vectors.

Results such as Abdalla and Zhivotovskiy \cite{abdalla2022} and Oliveira and
Rico \cite{oliveira2022} provide strong operator-norm guarantees under
adversarial contamination and $L_p$--$L_2$ marginal equivalence.  Minsker and
Wang \cite{minsker2022} study adversarial contamination and heavy tails with
adaptation to low-rank structure.  Their assumptions and conclusions should be
checked carefully against \eqref{eq:l4l2}; none should be invoked solely from
the absolute boundedness in \eqref{eq:bounded}.

Three defensible routes are available:
\begin{enumerate}[label=(\alph*)]
 \item assume a relative $L_4$--$L_2$ condition for the RFF distribution and
       apply an existing robust covariance theorem;
 \item exploit the special trigonometric structure of RFFs to prove an
       appropriate ridge-relative estimate directly;
 \item settle initially for an absolute operator-norm bound, then combine it
       with thresholding and \cref{eq:trunc-bound}.
\end{enumerate}
Route (c) is the most suitable first master's-level target.

\subsection{Estimating the cross-moment}

Covariance estimation alone is insufficient.  The normal equations also use
$c=\E[ZY]$.  Although $Z$ is bounded, $ZY$ inherits the tails of $Y$.  Under
bounded $Y$ it is a bounded-vector robust mean problem.  Under sub-Gaussian or
finite-moment noise it requires an estimator and theorem matched to those
assumptions.  An alternative is to use a robust regression method directly
rather than separately estimating $M$ and $c$; robust linear regression results
for Gaussian covariates illustrate what is possible in a more structured
setting \cite{diakonikolas2023}.

\section{A precise target theorem template}

The following is deliberately a template, not a claim.

\begin{question}[Finite-dimensional robust RFF ridge theorem]\label{q:main}
Under explicit assumptions on the conditional noise and the clean distribution
of $Z=\Phi(X)$, construct a polynomial-time estimator $\widetilde\beta$ such
that, with probability at least $1-\eta$,
\begin{equation}\label{eq:target-template}
 \norm{M^{1/2}(\widetilde\beta-\beta_\star)}_2
 \leq
 \cE_{\mathrm{sample}}(n,m,\eta)
 +\cE_{\mathrm{contam}}(\epsilon)
 +\cE_{\mathrm{ridge}}(\lambda)
 +R\tau^{r+1/2},
\end{equation}
where $\tau$ is a justified robust spectral-resolution scale and no positive
lower bound on $\lambda_{\min}(M)$ is assumed.
\end{question}

The theorem must state whether $\beta_\star$ is an unregularized population
coefficient, the best linear predictor in the RFF class, or the ridge target
$\beta_\lambda$.  These are different objects.  A clean first result should
compare with $\beta_\lambda$; approximation to $f_\star$ can then be handled in
a separate step.

\subsection{A candidate truncated estimator}

A concrete starting algorithm is:
\begin{enumerate}
 \item Compute robust estimates $\widehat M$ of $M$ and $\widehat c$ of $c$.
 \item Diagonalize $\widehat M=\sum_j\widehat\mu_j
       \widehat u_j\widehat u_j^\top$ and form
       $\widehat P_\tau=\sum_{\widehat\mu_j\geq\tau}
       \widehat u_j\widehat u_j^\top$.
 \item Return
 \begin{equation}\label{eq:truncated-estimator}
  \widetilde\beta_{\lambda,\tau}
  =\widehat P_\tau
  (\widehat P_\tau\widehat M\widehat P_\tau+\lambda\widehat P_\tau)^{\dagger}
  \widehat P_\tau\widehat c.
 \end{equation}
\end{enumerate}
Here $\dagger$ denotes the Moore--Penrose pseudoinverse on the retained
subspace.  The analysis should avoid requiring individual eigenvectors to be
identified when eigenvalues are clustered.  Projector or functional-calculus
arguments are preferable.

\section{From the finite-dimensional result back to KRR}

Let $\widehat f_{\mathrm{rob}}$ be the robust RFF estimator,
$f_{m,\lambda}$ the population or clean empirical RFF ridge comparator (the
choice must be fixed), and $f_\lambda$ the corresponding exact-kernel ridge
comparator.  Then
\begin{align}
 \norm{\widehat f_{\mathrm{rob}}-f_\star}_{L^2(P_X)}
 \leq{}&
 \underbrace{\norm{\widehat f_{\mathrm{rob}}-f_{m,\lambda}}_{L^2(P_X)}}_
 {\text{robust finite-dimensional estimation}}\nonumber\\
 &+\underbrace{\norm{f_{m,\lambda}-f_\lambda}_{L^2(P_X)}}_
 {\text{random-feature approximation}}\nonumber\\
 &+\underbrace{\norm{f_\lambda-f_\star}_{L^2(P_X)}}_
 {\text{clean ridge bias / statistical error}}.
 \label{eq:decomp}
\end{align}
This inequality is elementary, but each comparator must be defined on the same
population/sample probability space.  Mixing a population RFF solution with a
finite-sample exact KRR solution obscures which randomness each term contains.

Rahimi and Recht \cite{rahimi2007} establish the RFF construction.  Rudi and
Rosasco \cite{rudi2017} show that random-feature ridge regression can attain
optimal learning behavior with substantially fewer features than samples in
standard regimes.  Avron et al.\ \cite{avron2017} study spectral approximation
for KRR and explain the role of leverage-based feature sampling.  These works
provide tools for the middle term of \eqref{eq:decomp}, but the student must
match their regularization conventions and hypotheses to the robust theorem.

\section{The kernel-to-RFF source-condition question}

The finite-dimensional assumption $\beta=M^rv$ is mathematically convenient
but model-dependent: both $M$ and the coefficient vector depend on the realized
features.  The conceptually cleaner assumption is the kernel-level condition
$f_\star=T^rg$ in \eqref{eq:kernel-source}.

\begin{question}[Transfer of source regularity]
For which feature-sampling schemes, feature counts $m$, and regularization
levels $\lambda$ does the kernel source condition imply either
\begin{enumerate}[label=(\roman*)]
 \item a finite-dimensional source representation with controlled radius, or
 \item directly a small prediction contribution from the empirically
       unresolved RFF spectrum?
\end{enumerate}
\end{question}

Exact transfer may be too strong because the RFF subspace is random.  A more
realistic goal is an approximate statement with a residual term already
controlled by the RFF approximation error.  This question should be attempted
only after the finite-dimensional analysis is complete.

\section{Recommended project stages}

\begin{milestone}[Definitions and sanity checks; weeks 1--2]
Fix all target objects, distinguish covariance from second moment, derive the
normal equations and \eqref{eq:resolvent}, and reproduce
\cref{eq:trunc-bound}.  Construct a two-dimensional example showing why an
additive covariance estimate cannot safely whiten a tiny eigenvalue.
\end{milestone}

\begin{milestone}[Clean RFF baseline; weeks 3--4]
Implement exact KRR and RFF ridge regression for synthetic data.  Empirically
study how prediction error changes with $m$ and $\lambda$.  Read
\cite{rahimi2007,rudi2017,avron2017} and state one applicable approximation
result using the project's notation.
\end{milestone}

\begin{milestone}[Robust finite-dimensional baseline; weeks 5--7]
Begin with bounded responses and implement a robust mean estimator for $c$ and
a robust estimator for $M$.  Compare a plug-in ridge estimator, full whitening,
and spectral truncation under pair contamination.
\end{milestone}

\begin{milestone}[Perturbation theorem; weeks 8--10]
Prove a deterministic theorem: conditional on bounds for $\widehat M-M$ and
$\widehat c-c$, control prediction error of \eqref{eq:truncated-estimator} under
\eqref{eq:finite-source}.  Keep the stochastic robust-estimation result as a
separate corollary.
\end{milestone}

\begin{milestone}[Integration or counterexample; weeks 11--12]
Combine the deterministic result with one carefully verified robust estimation
theorem.  If the desired rate cannot follow from the available assumptions,
give a counterexample identifying the missing condition.  Either outcome is a
successful research contribution.
\end{milestone}

\begin{milestone}[Optional extension; weeks 13--14]
Investigate transfer from \eqref{eq:kernel-source}, leverage-weighted features,
or noise weaker than boundedness.  This is a stretch goal, not a prerequisite
for completing the project.
\end{milestone}

\section{Experiments}

Experiments should diagnose theory rather than merely compare test accuracy.
A useful synthetic design chooses a prescribed diagonal feature covariance,
places the target coefficient at controlled spectral locations, and corrupts
an $\epsilon$ fraction of complete pairs.  Record:
\begin{itemize}
 \item clean prediction error versus $\epsilon$, $m$, $n$, $\lambda$, and
       truncation threshold $\tau$;
 \item operator error $\norm{\widehat M-M}_\op$ and ridge-relative error
       $\Delta_\lambda$;
 \item the target energy $\sum_{\mu_j<\tau}\mu_j\beta_j^2$ below the threshold;
 \item sensitivity to an adversary concentrated in a low-variance direction;
 \item comparison with nonrobust RFF ridge, oracle clean RFF ridge, and exact
       clean KRR.
\end{itemize}
The oracle comparators are essential: they distinguish kernel approximation,
ordinary sampling error, and contamination error.

\section{Common pitfalls}

\begin{enumerate}
 \item \textbf{Bounded is not relatively hypercontractive.}  Do not infer
       \eqref{eq:l4l2} from \eqref{eq:bounded}.
 \item \textbf{Covariance is not the regression normal matrix.}  Account for
       the feature mean or use $M=\E[ZZ^\top]$ directly.
 \item \textbf{The response moment must be estimated.}  A robust covariance
       estimate does not by itself produce a regression estimator.
 \item \textbf{Ridge and truncation have different biases.}  Track both
       $\lambda$ and $\tau$ rather than merging them informally.
 \item \textbf{Eigenspaces need gaps; spectral functions often do not.}  Avoid
       claims about recovering individual eigenvectors unless a gap assumption
       is stated.
 \item \textbf{A displayed desired rate is not a theorem.}  Every final result
       must identify its estimator, probability, sample-size conditions,
       constants or parameter dependence, and computational status.
 \item \textbf{Condition on feature randomness consistently.}  Say whether a
       statement holds for most feature draws, for a fixed draw, or jointly
       over feature and sample randomness.
\end{enumerate}

\section{Concrete deliverables}

At the end of the project, the student should produce:
\begin{enumerate}
 \item a precise problem statement and literature review;
 \item a deterministic perturbation theorem for regularized or truncated
       moment-based regression;
 \item either a stochastic corollary under a verified robust-estimation result
       or a counterexample showing why stronger assumptions are necessary;
 \item reproducible synthetic experiments with oracle baselines;
 \item a clear account of what remains open, especially the transfer from a
       kernel source condition to finite random features.
\end{enumerate}

\section{Questions to settle with the advisor}

The following choices materially affect the theorem and should be fixed before
the student begins the main proof.
\begin{enumerate}
 \item Is the first theorem allowed to assume bounded responses, or should it
       begin immediately with conditionally sub-Gaussian noise?
 \item Should the primary comparator be the population RFF ridge solution
       $\beta_\lambda$ or the clean empirical RFF estimator?
 \item Is polynomial-time computation required, or is an information-theoretic
       estimator acceptable for the first result?
 \item Should complete-pair replacement remain the main model, or should the
       project first solve the simpler label-only contamination problem?
 \item Is the kernel-level source-condition transfer a required deliverable or
       explicitly a stretch goal?
\end{enumerate}

\section{Reading guide}

The following order is reasonable for a master's student.
\begin{enumerate}
 \item Read Rahimi and Recht \cite{rahimi2007} for Bochner's theorem and the
       basic RFF construction.
 \item Read selected sections of Rudi and Rosasco \cite{rudi2017} for learning
       with random features, then Avron et al.\ \cite{avron2017} for the
       regularized spectral-approximation viewpoint.
 \item Review KRR source conditions and effective dimension in Caponnetto and
       De Vito \cite{caponnetto2007}, together with Sutherland's correction
       \cite{sutherland2017}.
 \item Read a robust statistics overview such as Diakonikolas and Kane
       \cite{diakonikolas2019} to learn the strong contamination model and
       filtering paradigm.
 \item Only then study robust covariance papers
       \cite{abdalla2022,oliveira2022,minsker2022}, writing down precisely which
       moment assumption each theorem uses.
\end{enumerate}

\begin{thebibliography}{99}

\bibitem{abdalla2022}
P. Abdalla and N. Zhivotovskiy.
\newblock Covariance estimation: optimal dimension-free guarantees for
adversarial corruption and heavy tails.
\newblock \emph{arXiv:2205.08494}, 2022.
\newblock \url{https://arxiv.org/abs/2205.08494}.

\bibitem{avron2017}
H. Avron, M. Kapralov, C. Musco, C. Musco, A. Velingker, and A. Zandieh.
\newblock Random Fourier features for kernel ridge regression: approximation
bounds and statistical guarantees.
\newblock In \emph{Proceedings of the 34th International Conference on Machine
Learning}, 2017.  Extended version: \url{https://arxiv.org/abs/1804.09893}.

\bibitem{caponnetto2007}
A. Caponnetto and E. De Vito.
\newblock Optimal rates for the regularized least-squares algorithm.
\newblock \emph{Foundations of Computational Mathematics}, 7:331--368, 2007.

\bibitem{diakonikolas2019}
I. Diakonikolas and D. M. Kane.
\newblock Recent advances in algorithmic high-dimensional robust statistics.
\newblock \emph{arXiv:1911.05911}, 2019.
\newblock \url{https://arxiv.org/abs/1911.05911}.

\bibitem{diakonikolas2023}
I. Diakonikolas, D. M. Kane, A. Pensia, and T. Pittas.
\newblock Near-optimal algorithms for Gaussians with Huber contamination: mean
estimation and linear regression.
\newblock \emph{arXiv:2312.01547}, 2023.
\newblock \url{https://arxiv.org/abs/2312.01547}.

\bibitem{minsker2022}
S. Minsker and L. Wang.
\newblock Robust estimation of covariance matrices: adversarial contamination
and beyond.
\newblock \emph{arXiv:2203.02880}, 2022.
\newblock \url{https://arxiv.org/abs/2203.02880}.

\bibitem{oliveira2022}
R. I. Oliveira and Z. F. Rico.
\newblock Improved covariance estimation: optimal robustness and sub-Gaussian
guarantees under heavy tails.
\newblock \emph{arXiv:2209.13485}, 2022.
\newblock \url{https://arxiv.org/abs/2209.13485}.

\bibitem{rahimi2007}
A. Rahimi and B. Recht.
\newblock Random features for large-scale kernel machines.
\newblock In \emph{Advances in Neural Information Processing Systems 20}, 2007.
\newblock \url{https://people.eecs.berkeley.edu/~brecht/papers/07.rah.rec.nips.pdf}.

\bibitem{rudi2017}
A. Rudi and L. Rosasco.
\newblock Generalization properties of learning with random features.
\newblock In \emph{Advances in Neural Information Processing Systems 30}, 2017.
\newblock \url{https://papers.nips.cc/paper/2017/hash/61b1fb3f59e28c67f3925f3c79be81a1-Abstract.html}.

\bibitem{sutherland2017}
D. J. Sutherland.
\newblock Fixing an error in Caponnetto and De Vito (2007).
\newblock \emph{arXiv:1702.02982}, 2017.
\newblock \url{https://arxiv.org/abs/1702.02982}.

\end{thebibliography}

\end{document}
